\section{Preliminaries}
\label{sec:preliminaries}

This section formalizes the LLM routing problem as a cost-aware contextual bandit, introduces the feature representation, and reviews the two algorithmic building blocks---LinUCB and Corralling---that underpin the banditGPT architecture.

\subsection{Problem Formulation}
\label{sec:formulation}

Let $x_t \in \mathbb{R}^d$ denote the context vector (embedding) for a user prompt at time $t$. We maintain a dynamic registry of models (arms) $\mathcal{M}_t$. For each request, the policy $\pi$ selects an arm $a_t \in \mathcal{M}_t$, observing a scalar reward $r_t \in [0, 1]$ (quality) and incurring a deterministic cost $c_t \in \mathbb{R}^+$ (inference price).

Unlike standard regret minimization, our objective is to maximize quality subject to a budget constraint, formulated as a Lagrangian relaxation with a user-defined penalty $\lambda$:
\begin{equation}
\text{Maximize } \sum_{t=1}^T \left( \mathbb{E}[r_t(a_t) | x_t] - \lambda \cdot c_t(a_t) \right)
\end{equation}
where $\lambda$ controls the trade-off between performance and economy. A high $\lambda$ forces the router to identify ``easy'' prompts where cheaper models suffice, while a low $\lambda$ prioritizes quality regardless of cost.

\subsection{Background: LinUCB and Expert Corralling}
\label{sec:background}

\paragraph{LinUCB.}
LinUCB~\cite{li2010contextual} models the expected reward of each arm $a$ as a linear function of the context: $\mathbb{E}[r | x] = x^\top \theta^*_a$. It maintains a regularized least-squares estimate $\hat{\thetavec}_a$ and an associated confidence set, selecting the arm that maximizes an upper confidence bound on the composite utility. LinUCB provides a natural framework for contextual routing because per-prompt embeddings allow the policy to vary model selection with input content.

\paragraph{Expert Corralling.}
Corralling~\cite{agarwal2017corralling} is a meta-learning strategy that maintains a portfolio of expert policies, dynamically allocating probability mass among them via exponential weight updates. By ensuring non-zero probability on each expert, it provides asymptotic recovery guarantees when the optimal expert changes over time. In our setting, Corralling enables safe hedging against prior misspecification by maintaining both a warm-started and a tabula rasa expert.

\subsection{Unified Feature Space and Empirical Validation}
\label{sec:motivation}

A contextual bandit is only useful if prompt features carry actionable routing signal. We embed prompts using Sentence-BERT~\cite{reimers2019sentence} (\path{all-MiniLM-L6-v2}, 384D) and reduce to 32 dimensions using unsupervised PCA trained on 80k routing prompts from RouteLLM~\cite{ong2024routellm}.

To validate that these features predict model preference, we evaluate on 750 held-out prompts comparing a frontier model (GPT-4-Turbo) against a cheaper open-weights alternative (Mixtral-8x7B). LLM routing literature often assumes \emph{monotonicity}: that allocating more budget to a frontier model consistently yields equal or higher quality~\cite{ong2024routellm, chen2023frugalgpt}. Our held-out data challenge this assumption: human judges preferred the cheaper model on 14.1\% of prompts, while the frontier model won on only 13.1\%, with 72.8\% ties (Table~\ref{tab:outcome_distribution}).

\begin{table}[htbp]
\centering
\small
\caption{Outcome distribution on held-out data ($N=750$).}
\label{tab:outcome_distribution}
\begin{tabular}{@{}lrr@{}}
\toprule
\textbf{Outcome} & \textbf{Count} & \textbf{Proportion} \\
\midrule
Frontier model wins & 98 & 13.1\% \\
Tie & 546 & 72.8\% \\
Cheaper model wins & 106 & 14.1\% \\
\bottomrule
\end{tabular}
\end{table}

Crucially, this preference heterogeneity is \emph{predictable} from the embedding. When prompts are projected onto the first principal component (PC1) of our PCA feature space, we observe a significant Spearman correlation with the reward gap ($R_{\text{Frontier}} - R_{\text{Cheaper}}$): $\rho = -0.370$, $p < 0.0001$ (Figure~\ref{fig:lmsys_holdout_structure}A). This signal exceeds all 100 random orthonormal projections (Figure~\ref{fig:lmsys_holdout_structure}A annotation), and the routing opportunity grows monotonically across PC1 quintiles (Figure~\ref{fig:lmsys_holdout_structure}B), confirming that the PCA axis separates feature regions with distinct model preferences.

% Figure 1 Caption: Contextual Sensitivity across K=10 Models

\begin{figure*}[t]
\centering
\includegraphics[width=0.85\textwidth]{figures/figure1_k10_contextual.png}
\caption{\textbf{Model quality depends on prompt context, with heterogeneous sensitivity across $K{=}10$ models ($N{=}750$, holdout).}
\textbf{(A)}~OLS regression of mean reward ($\text{mean}(v \times c)$, from a 3--4 judge panel) on standardised PC1 for each model, with 95\% confidence bands.
Lines fan and cross: GPT-4-Turbo declines steeply ($\gamma{=}{-}0.254$) while GPT-4.1 is nearly flat ($\gamma{=}{-}0.032$), so no single model dominates the full feature space.
A likelihood-ratio test rejects the shared-slope null ($\chi^2{=}771.3$, $\text{df}{=}54$, $p < 10^{-15}$; six PCA components).
\textbf{(B)}~Forest plot of per-model contextual slopes $\gamma_m$ with bootstrap 95\% CIs ($10\,000$ resamples).
All ten CIs exclude zero, confirming that every model's quality varies with prompt features---but at rates differing by nearly an order of magnitude.
PCA trained on 80K RouteLLM battle prompts (independent dataset); reward computed as $\text{mean}(\text{vote} \times \text{confidence})$ over the judge panel; PC1 is one of 32 components used by the router.}
\Description{Two-panel figure. Panel~A shows ten OLS regression lines of reward versus standardised PC1, with shaded 95\% confidence bands; lines separate and cross, indicating different models are optimal in different feature regions. Panel~B is a forest plot of per-model contextual slopes with bootstrap confidence intervals; all ten intervals exclude zero, and magnitudes range from 0.032 to 0.254.}
\label{fig:contextual_heterogeneity}
\end{figure*}


These results confirm that our feature space captures routing-relevant signals suitable for online adaptation: embeddings carry actionable model-preference information that a contextual bandit can exploit to route each prompt to the most appropriate model.

